\section{Bounded Synthesis of the K-Level Model Core}

\label{sec:oc-core-1-3-unified-synthesis}

This chapter replaces the historical generated synthesis status sheet with a publication-grade exposition. The scientific task is not to announce that every domain has been completed. The task is to show how the K-level witness taxonomy, typed carriers, realization predicates, finite witnesses, proof sheets, and bounded replay rows fit together as one model-core route.


The chapter should be read with three limits in mind. First, a K-level distinction is promoted only where a witness or theorem route makes the distinction do work. Second, numerical rows are reconstruction and target-blind replay evidence inside declared protocols, not a claim that the whole scientific domain has been proved. Third, complete scientific coverage and unrestricted cross-science comparison are outside the promoted 1.3.3 public claim surface.


\subsection{How the synthesis is organized}

The synthesis follows a reader-facing order rather than a machine-log order. Each level states its role, its current evidence class, the kind of reviewer challenge it answers, and the condition that would reopen the claim. Exhaustive hashes, verdict fields, and case rows remain in the evidence package and are cited from the monolith rather than pasted into the main prose.


\subsection{K0: proto-continuum and zero-continuumness boundary}

K0 contributes proto-continuum and zero-continuumness boundary. Its public release evidence is the K-zero boundary theorem, the K0 resolution theorem, and finite semantic witnesses. The reviewer challenge is whether the model can distinguish absence of live realization from low-level structure. The release claim is therefore scoped: the level is part of the OC model-core grammar where its declared witness obligations are met.


For K0, the claim reopens if a lower-level projection preserves every declared witness, if the cited proof sheet loses an assumption, if the finite semantic case can be tampered with without failing, or if a domain row no longer reproduces its stated observable under the pinned protocol. This is the evidence-inclusion rule for this level: it must explain what it adds, how the evidence binds it, and what would make the wording too strong.


\subsection{K1: minimal separability and first structured distinction}

K1 contributes minimal separability and first structured distinction. Its public release evidence is typed carrier definitions and boundary predicates. The reviewer challenge is whether a lower projection can preserve the declared distinction. The release claim is therefore scoped: the level is part of the OC model-core grammar where its declared witness obligations are met.


For K1, the claim reopens if a lower-level projection preserves every declared witness, if the cited proof sheet loses an assumption, if the finite semantic case can be tampered with without failing, or if a domain row no longer reproduces its stated observable under the pinned protocol. This is the evidence-inclusion rule for this level: it must explain what it adds, how the evidence binds it, and what would make the wording too strong.


\subsection{K2: early organization above separability}

K2 contributes early organization above separability. Its public release evidence is operator semantics and low-level finite checks. The reviewer challenge is whether update rules add anything beyond static classification. The release claim is therefore scoped: the level is part of the OC model-core grammar where its declared witness obligations are met.


For K2, the claim reopens if a lower-level projection preserves every declared witness, if the cited proof sheet loses an assumption, if the finite semantic case can be tampered with without failing, or if a domain row no longer reproduces its stated observable under the pinned protocol. This is the evidence-inclusion rule for this level: it must explain what it adds, how the evidence binds it, and what would make the wording too strong.


\subsection{K3: organized physical structures}

K3 contributes organized physical structures. Its public release evidence is bounded replay rows and proof-boundary conditions. The reviewer challenge is whether physical examples are merely relabeled standard quantities. The release claim is therefore scoped: the level is part of the OC model-core grammar where its declared witness obligations are met.


For K3, the claim reopens if a lower-level projection preserves every declared witness, if the cited proof sheet loses an assumption, if the finite semantic case can be tampered with without failing, or if a domain row no longer reproduces its stated observable under the pinned protocol. This is the evidence-inclusion rule for this level: it must explain what it adds, how the evidence binds it, and what would make the wording too strong.


\subsection{K4: chemical binding and composition routes}

K4 contributes chemical binding and composition routes. Its public release evidence is chemistry replay rows, comparator notes, and operator boundaries. The reviewer challenge is whether composition is explained or just renamed. The release claim is therefore scoped: the level is part of the OC model-core grammar where its declared witness obligations are met.


For K4, the claim reopens if a lower-level projection preserves every declared witness, if the cited proof sheet loses an assumption, if the finite semantic case can be tampered with without failing, or if a domain row no longer reproduces its stated observable under the pinned protocol. This is the evidence-inclusion rule for this level: it must explain what it adds, how the evidence binds it, and what would make the wording too strong.


\subsection{K5: excitable and biological organization}

K5 contributes excitable and biological organization. Its public release evidence is biology replay rows and lifecycle semantics. The reviewer challenge is whether liveness, residue, and rebirth are formally separable. The release claim is therefore scoped: the level is part of the OC model-core grammar where its declared witness obligations are met.


For K5, the claim reopens if a lower-level projection preserves every declared witness, if the cited proof sheet loses an assumption, if the finite semantic case can be tampered with without failing, or if a domain row no longer reproduces its stated observable under the pinned protocol. This is the evidence-inclusion rule for this level: it must explain what it adds, how the evidence binds it, and what would make the wording too strong.


\subsection{K6: cognitive and agent-scale organization}

K6 contributes cognitive and agent-scale organization. Its public release evidence is phenomenon model cards and identity-boundary rules. The reviewer challenge is whether identity claims preserve declared invariants. The release claim is therefore scoped: the level is part of the OC model-core grammar where its declared witness obligations are met.


For K6, the claim reopens if a lower-level projection preserves every declared witness, if the cited proof sheet loses an assumption, if the finite semantic case can be tampered with without failing, or if a domain row no longer reproduces its stated observable under the pinned protocol. This is the evidence-inclusion rule for this level: it must explain what it adds, how the evidence binds it, and what would make the wording too strong.


\subsection{K7: institutional and social organization}

K7 contributes institutional and social organization. Its public release evidence is systems model cards and falsifier routes. The reviewer challenge is whether social examples become metaphor rather than model instances. The release claim is therefore scoped: the level is part of the OC model-core grammar where its declared witness obligations are met.


For K7, the claim reopens if a lower-level projection preserves every declared witness, if the cited proof sheet loses an assumption, if the finite semantic case can be tampered with without failing, or if a domain row no longer reproduces its stated observable under the pinned protocol. This is the evidence-inclusion rule for this level: it must explain what it adds, how the evidence binds it, and what would make the wording too strong.


\subsection{K8: civilizational and multi-layer coordination}

K8 contributes civilizational and multi-layer coordination. Its public release evidence is systems/civilizational replay rows and comparator boundaries. The reviewer challenge is whether aggregate patterns remain tied to observable protocols. The release claim is therefore scoped: the level is part of the OC model-core grammar where its declared witness obligations are met.


For K8, the claim reopens if a lower-level projection preserves every declared witness, if the cited proof sheet loses an assumption, if the finite semantic case can be tampered with without failing, or if a domain row no longer reproduces its stated observable under the pinned protocol. This is the evidence-inclusion rule for this level: it must explain what it adds, how the evidence binds it, and what would make the wording too strong.


\subsection{K9: theory-level self-description}

K9 contributes theory-level self-description. Its public release evidence is claim/evidence governance and proof-register links. The reviewer challenge is whether the release can distinguish public claims from research targets. The release claim is therefore scoped: the level is part of the OC model-core grammar where its declared witness obligations are met.


For K9, the claim reopens if a lower-level projection preserves every declared witness, if the cited proof sheet loses an assumption, if the finite semantic case can be tampered with without failing, or if a domain row no longer reproduces its stated observable under the pinned protocol. This is the evidence-inclusion rule for this level: it must explain what it adds, how the evidence binds it, and what would make the wording too strong.


\subsection{K10: meta-theoretic comparison}

K10 contributes meta-theoretic comparison. Its public release evidence is novelty register and source-backed comparator matrix. The reviewer challenge is whether the work is only a reframing of prior art. The release claim is therefore scoped: the level is part of the OC model-core grammar where its declared witness obligations are met.


For K10, the claim reopens if a lower-level projection preserves every declared witness, if the cited proof sheet loses an assumption, if the finite semantic case can be tampered with without failing, or if a domain row no longer reproduces its stated observable under the pinned protocol. This is the evidence-inclusion rule for this level: it must explain what it adds, how the evidence binds it, and what would make the wording too strong.


\subsection{K11: methodological and publication quality}

K11 contributes methodological and publication quality. Its public release evidence is editorial criteria, science-state summaries, and reproducibility checks. The reviewer challenge is whether the artifact pipeline can prevent unsupported publication surfaces. The release claim is therefore scoped: the level is part of the OC model-core grammar where its declared witness obligations are met.


For K11, the claim reopens if a lower-level projection preserves every declared witness, if the cited proof sheet loses an assumption, if the finite semantic case can be tampered with without failing, or if a domain row no longer reproduces its stated observable under the pinned protocol. This is the evidence-inclusion rule for this level: it must explain what it adds, how the evidence binds it, and what would make the wording too strong.


\subsection{K12: cross-model coherence under declared assumptions}

K12 contributes cross-model coherence under declared assumptions. Its public release evidence is bounded synthesis evidence and adversarial review. The reviewer challenge is whether the declared model grammar remains coherent when projected across domains. The release claim is therefore scoped: the level is part of the OC model-core grammar where its declared witness obligations are met.


For K12, the claim reopens if a lower-level projection preserves every declared witness, if the cited proof sheet loses an assumption, if the finite semantic case can be tampered with without failing, or if a domain row no longer reproduces its stated observable under the pinned protocol. This is the evidence-inclusion rule for this level: it must explain what it adds, how the evidence binds it, and what would make the wording too strong.


\subsection{What the bounded synthesis licenses}

The bounded synthesis licenses model-core claims about typed continuants, lawful realization, liveness, residue, boundary classification, hybrid update/flow structure, K-level witness distinctions, and selected target-blind replay rows. It does not license a public statement that all phenomena in the declared release domains are already numerically predicted, nor a statement that all prior science has been surpassed.


This distinction is a release requirement rather than a rhetorical caution. The public manuscript must teach the synthesis, cite the evidence, and preserve the ambition, while keeping the broader research obligations visible as obligations. That is the standard used by the 1.3.3 editorial and release criteria.
